% drafts/10-concerns.tex — Addressing Potential Concerns
%
% Target length: ~2 pages (two-column PRA)
% Proactive Q&A addressing likely reviewer and reader questions.

We anticipate several natural questions about the scope, proof
technique, and significance of the results in this paper.  Rather than
leave these for reviewer queries, we address them directly.

%%====================================================================
\subsection{Why perform exhaustive enumeration?}
\label{sec:concern-proof}

Our proof of the star-tree theorem (\cref{thm:star-tree}) is
structural, identifying a generic failure mode for depth-$\ge 2$ trees.
Why then do we also present exhaustive enumeration results?

The enumeration serves as an independent, model-free validation.
Structural proofs rely on identifying specific mechanisms (here, the
overlap of remainder sets along a path), which might have subtle
exceptions or edge cases.  By blindly generating all $n^{n-1}$ trees
for small $n$ and symbolically verifying every anticommutator, we
confirm that:
\begin{enumerate}
  \item No unexpected valid trees exist (e.g., highly symmetric
    non-star graphs that might bypass the failure mode).
  \item The failure is indeed universal for all non-star structures.
  \item The count of valid trees corresponds exactly to the number
    of stars ($n$), with no "accidental" successes.
\end{enumerate}

This combination of analytical proof and computational verification
provides the highest possible confidence in the result.

%%====================================================================
\subsection{Could a different index-set definition work?}
\label{sec:concern-analytical}

The failure of Construction~A suggests that the specific definitions
of update, parity, and remainder sets in the Seeley--Richard--Love (SRL)
framework are responsible.  Could one simply redefine $\Rset{j}$ to fix
the problem?

Almost certainly.  The Bravyi--Kitaev encoding uses the Fenwick tree
(depth $\sim \log n$) but requires a different set of index
definitions (Construction~F).  It is likely that other trees admit
valid index sets if derived specifically for their structure.
However, our result shows that the \emph{generic} SRL algorithm---which
claims to work for any tree---fails for all but stars.  The path-based
Construction~B solves this by abandoning index sets entirely in favor
of explicit string generation.

%%====================================================================
\subsection{Does the star-tree result generalise?}
\label{sec:concern-generalise}

One might ask whether the star-tree limitation is specific to the SRL
index-set construction or reflects a deeper algebraic constraint.

The answer is the former:\ it is specific to Construction~A.  The
path-based Construction~B, using the same Majorana intermediate
representation, works for \emph{all} trees --- including non-star,
non-monotonic trees of arbitrary depth.  The star-tree theorem is thus
a result about the \emph{index-set derivation procedure}, not about
tree-based encodings in general.

This matters because it clarifies the relationship between the two
frameworks:
\begin{itemize}
  \item Index-set schemes are more compact to specify (three
    functions) but structurally restricted in which trees they can
    handle.
  \item Tree encodings are more general (universal) but require
    explicit tree traversal.
  \item The two frameworks agree on their overlap domain (star
    trees / JW / Parity) but diverge for all richer topologies.
\end{itemize}

%%====================================================================
\subsection{Is the bijective proof of $|M(n)| = (n{-}1)!$ known?}
\label{sec:concern-bijective}

The identity $|M(n)| = (n{-}1)!$ (\cref{prop:monotonic-count}) is
verified computationally for $n \le 6$.  We are not aware of this
specific identity appearing in the encoding literature, though the
connection to heap-ordered trees is well-studied in
enumerative combinatorics~\cite{bergeron1992}.

A monotonic labelling (parent $>$ children everywhere) is equivalent
to a \emph{heap ordering} on a rooted tree.  The number of heap
orderings summed over all unlabelled rooted trees on $n$ vertices is
known to be $(n{-}1)!$ (this follows from the ``parking function''
bijection or from the recursive decomposition of increasing
trees~\cite{bergeron1992}).  We believe our computational result
is a rediscovery of this classical identity in the encoding context,
but establishing the bijection explicitly in encoding-theoretic
terms would be a useful pedagogical contribution.

%%====================================================================
\subsection{Why not claim computational speedup?}
\label{sec:concern-speedup}

The encoding pipeline has $O(n^5)$ complexity for molecular
Hamiltonians, identical to OpenFermion~\cite{mcclean2020} and
PennyLane~\cite{bergholm2022}.  We deliberately do not claim
computational improvement, because there is none.

The advantage is \emph{structural}, not asymptotic:
\begin{itemize}
  \item \textbf{Correctness by construction:} the type system
    (\cref{app:datatypes}) makes certain classes of bugs
    impossible (phase errors, unmatched operators, forgotten
    cancellations).
  \item \textbf{Extensibility:} adding a new encoding requires
    3--5 lines (a tree or three set functions), not hundreds.
  \item \textbf{Exact phases:} symbolic $\Zi$ tracking eliminates
    floating-point phase accumulation (\cref{thm:exactness}).
  \item \textbf{Unified pipeline:} fermionic, bosonic, and mixed
    systems share the same infrastructure.
\end{itemize}

We believe these are valuable properties for a quantum chemistry
middleware layer, but they are engineering contributions, not
algorithmic breakthroughs.  We are careful to distinguish the two.

%%====================================================================
\subsection{How does this compare to PennyLane and OpenFermion?}
\label{sec:concern-comparison}

A fair comparison must distinguish \emph{scope} from
\emph{capability}:

\begin{center}
\begin{tabular}{@{}p{2.5cm}p{2cm}p{2cm}p{2cm}@{}}
\toprule
& \textbf{FockMap} & \textbf{OpenFermion} & \textbf{PennyLane} \\
\midrule
Encoding as data structure &
  \checkmark & --- & --- \\
Custom tree encodings &
  \checkmark & --- & --- \\
Exact phase tracking &
  $\Zi$ symbolic & float & float \\
Bosonic operators &
  \checkmark & \checkmark & \checkmark \\
Mixed normal ordering &
  \checkmark & --- & --- \\
Circuit synthesis &
  --- & \checkmark & \checkmark \\
VQE / optimization &
  --- & --- & \checkmark \\
\bottomrule
\end{tabular}
\end{center}

FockMap occupies a different niche: it is a \emph{Hamiltonian
construction} library, not a quantum circuit or variational framework.
It produces the Pauli Hamiltonian; downstream tools (Qiskit, Cirq,
PennyLane) consume it for circuit synthesis and execution.  The
comparison is complementary, not competitive.

%%====================================================================
\subsection{Is the branching factor $\le 3$ restriction fundamental?}
\label{sec:concern-branching}

Construction~B (path-based encoding) requires each node to have at
most three descending links, labelled $X$, $Y$, $Z$.  This
restriction arises because there are exactly three non-identity
single-qubit Pauli operators.

A natural question is whether higher-arity trees could be supported by
using multi-qubit Pauli labels on links.  The answer is yes in
principle: one could label each link with a multi-qubit Pauli
substring and allocate multiple qubits per tree node.  However, this
would change the qubit-to-mode ratio ($n$ qubits for $n$ modes is no
longer guaranteed) and complicate the weight analysis.  Whether such
generalized encodings offer any practical advantage is an open
question.

For the standard one-qubit-per-mode setting, the branching factor~3
restriction implies that the balanced ternary tree is the maximum-arity
balanced tree, and its $O(\log_3 n)$ weight is optimal within this
class.
